\section{Assembling the original Galois representation of a definite unitary form}
\label{sec:galois-assembly-v44}

\status{This chapter reconstructs \texttt{GalUnitaryT} and its
\texttt{GuWrap} consumer.  It uses the original polarized cuspidal
attachment input \texttt{RowT}, the original local representation and
period foundations, and the base-change/Speh construction BCS18--BCS19.
These inputs are identified separately from the proofs below.  The
checked published generic-unitary input is BCS8; GUA4 supplies its
needed elementary ordering consequence.  We preserve every positive rank,
every Speh length, every odd coefficient prime including $3$, and
the universal quantifiers over local correspondences and period
bridges.  In particular, the local conclusion retains the source's
\texttt{Gd} condition.  This is a reconstruction of an installation
input, not an unconditional proof of RH or GRH.}

Let $E$ be a CM number field, $E^+$ its maximal totally real
subfield, and $\Omega/E$ the specified algebraic closure.  Suppose
$E/E^+$ is unramified at every finite place.  Fix $n\geq1$, a prime
$\ell\ne2$, and the given field isomorphism
$\iota:\overline{\mathbb Q}_\ell\simeq\mathbb C$.
Let $\sigma$ be the original irreducible automorphic subspace of
functions on the definite unitary adelic group $U(\mathbb A_{E^+})$,
with the right-translation action; this is \texttt{IsAutRepU}.
No continuity of $\iota$ for the complex topology is assumed.

Write $\nu=|\det|_F$ over a finite completion $F=E_w$, with
$|\varpi|_F=q_w^{-1}$.  The source conventions are
\[
 \operatorname{algNorm}(\tau)=\tau\nu^{(1-\dim\tau)/2},\qquad
 x(t)=x.\texttt{lpTw}(t),\qquad
 \epsilon(\operatorname{Frob}^{\rm ar}_w)=q_w.        \tag{GUA.1}
\]
The twist $x(t)$ multiplies the Weil action on an arithmetic
Frobenius by $q_w^t$; it leaves monodromy unchanged.  Thus its
geometric Frobenius eigenvalues are multiplied by $q_w^{-t}$.
All real powers of positive numbers in this chapter are the
positive real powers, embedded in $\mathbb C$.

For a Weil--Deligne isomorphism class $y=(r_y,N_y)$, put
\[
 \operatorname{Gd}(y)\quad\Longleftrightarrow\quad
 N_y=0\quad\hbox{and}\quad
 \det(X-r_y([u,v]))=(X-1)^{\dim y}
                       \quad(u,v\in W_{E_w}).       \tag{GUA.2}
\]
This is the literal \texttt{Tor}$\wedge$\texttt{mono0} predicate.
It is not being strengthened to a specified diagonalization.
For every valuation $A$ of $\Omega$ over $w$ and every map
$\mathrm{rc}:\operatorname{Irr}\operatorname{GL}(E_w)
\longrightarrow\operatorname{WD}(A)$ satisfying the full original
\texttt{IsRecT}, the rank-one law uses that correspondence's own
local Artin map $a:W_A\to E_w^\times$.  Its degree convention is
$\operatorname{ord}_w(a(g))=-\deg(g)$.

We use the following exact compatibility condition.  A triple
$(R,p,H)$ satisfies \texttt{CompatG} if $R$ is continuous and
semisimple and the following statements hold for every $w$.
First, for every $A,h,\mathrm{rc}$ as above, at $w\nmid\ell$
there is an actual Weil--Deligne parameter $x$ of $R|_{D_A}$
such that
\[
             \operatorname{Gd}(p_{w,A,\mathrm{rc}})
                    \ \Longrightarrow\ x=p_{w,A,\mathrm{rc}}.
                                                        \tag{GUA.3}
\]
Second, for every original coefficient ring $O$, local bridge
$B_w$ and witness \texttt{Induces}, goodness of that same
parameter implies potential semistability in the specified
period model $\texttt{Loc.PP}(\ell)$ and the exact equality
\[
 \operatorname{wdAt}(R|_{D_A},B_w,\iota)
                              =p_{w,A,\mathrm{rc}}.       \tag{GUA.4}
\]
Finally, if the parameters are good for \emph{all}
$A,h,\mathrm{rc}$ at $w$, the representation is de Rham for
every such bridge, and its filtration jumps at every induced
embedding are the multiset $H_\tau$; this is
\texttt{HasHTAt}.  The quantification includes every finite
coefficient extension and every inducing embedding appearing
in that definition.  No existence of a bridge at an unrelated
residue prime is asserted.

\subsection{Local parameter identities for every permitted correspondence}

% BEGIN GUA1
\paragraph{GUA1. Goodness under finite sums and unramified twists.}
\label{claim:GUA1}
For Weil--Deligne classes $x,y$ and $t\in\mathbb R$,
\[
 \operatorname{Gd}(x\oplus y)
       \Longleftrightarrow\operatorname{Gd}(x)\wedge\operatorname{Gd}(y),
 \qquad
 \operatorname{Gd}(x(t))\Longleftrightarrow\operatorname{Gd}(x).
\]
The corresponding statement holds for every finite list,
including the empty list.  An unramified class is good.
\begin{proof}
Choose representatives.  Monodromy on a sum is the block
matrix $N_x\oplus N_y$, which vanishes exactly when both
blocks vanish.  For each commutator the characteristic
polynomial is the product of the two block characteristic
polynomials.  A product of two monic complex polynomials is
$(X-1)^{a+b}$ exactly when its degree-$a$ and degree-$b$
factors are $(X-1)^a$ and $(X-1)^b$: all their complex roots
must be $1$.  This proves the first equivalence.  An
unramified twist factors through the degree map
$W_A\to\mathbb Z$, so it is $1$ on every commutator and
does not change $N$.  Induction proves the list statement;
the zero class satisfies both conditions.  Finally a
commutator has degree zero and hence lies in inertia.  An
unramified class is trivial there and has zero monodromy.
These arguments are invariant under change of representative.
\end{proof}
% END GUA1

% BEGIN GUA2
\paragraph{GUA2. The full \texttt{PlainGd} statement.}
\label{claim:GUA2}
For every $A,h,\mathrm{rc}$ satisfying \texttt{IsRecT}, every
rank $a\geq0$ and every smooth irreducible $\tau$ of
$\operatorname{GL}_a(E_w)$,
\[
 \texttt{PlainTor}(\operatorname{cl}\tau)
            \Longrightarrow
       \operatorname{Gd}\bigl(\mathrm{rc}(\operatorname{algNorm}\tau)\bigr).
                                                        \tag{GUA.5}
\]
\begin{proof}
The original \texttt{PlainTor} says that the cuspidal support
consists of rank-one characters $\xi_1,\ldots,\xi_a$ and is
unlinked.  Thus the singleton segments are pairwise unlinked;
the retained Bernstein--Zelevinsky irreducibility criterion
makes their normalized principal series irreducible.  The
support condition makes $\tau$ a subquotient of this principal
series, hence isomorphic to it.  Order the singleton segments
so none precedes a later one.  The standard-cosocle theorem
then identifies this same representation with their Langlands
quotient.  There is no requirement that the characters be
distinct: repeated singletons do not form linked segments.

Write $a_{\mathrm{rc}}$ for the local Artin map contained in
the given \texttt{IsRecT}.  Its rank-one and Langlands laws
give the equality of actual isomorphism classes
\[
 \mathrm{rc}(\operatorname{cl}\tau)
       =\bigoplus_{j=1}^a (\xi_j\circ a_{\mathrm{rc}},0).
\]
Each character is trivial on commutators and has zero
monodromy.  GUA1 proves goodness.  The real-twist law of
that same correspondence identifies the parameter of
$\operatorname{algNorm}\tau$ with the twist by $(1-a)/2$;
GUA1 applies once more.  Empty induction and the dimension
law give the zero parameter when $a=0$.  The construction
used no choice of a preferred correspondence or Artin map,
so it proves the stated universal quantifier.
\end{proof}
% END GUA2

% BEGIN GUA3
\paragraph{GUA3. Irreducible normalized products have summed parameters.}
\label{claim:GUA3}
If an irreducible smooth representation $\tau$ is isomorphic
to a finite normalized parabolic product of irreducibles
$\tau_i$, then every \texttt{IsRecT} satisfies
\[
          \mathrm{rc}(\operatorname{cl}\tau)
                      =\bigoplus_i\mathrm{rc}(\operatorname{cl}\tau_i).
                                                        \tag{GUA.6}
\]
In particular this proves the literal \texttt{ProdT}, with
the original block ordering and dimension equality.
\begin{proof}
For two factors choose their Langlands multisegments
$\mathfrak m,\mathfrak n$ using the retained exhaustion
theorem.  Order $\mathfrak m+\mathfrak n$ and form its
standard module.  The retained standard-cosocle theorem
gives its irreducible quotient $L(\mathfrak m+\mathfrak n)$.
The precise multisegment product theorem places this quotient
as a subquotient of $L(\mathfrak m)\times L(\mathfrak n)$.
If the latter is irreducible the two are isomorphic.
Now apply the Langlands law in \texttt{IsRecT}.  Its
\texttt{segSum} attaches the segment parameter to each
member of the list and takes their direct sum.  Concatenating
the two lists concatenates these blocks, so it gives (GUA.6).
For several factors, the same product theorem repeatedly
puts the Langlands quotient of the total multisegment into
the full induction; irreducibility of that full induction
again identifies it with $\tau$.  This argument does not
require an unstated irreducibility of an intermediate product.
The empty case has dimension zero and the unique zero
Weil--Deligne class.  Reassociation and the source's
$\texttt{stdBlocks}$ coordinate permutation give an actual
intertwining isomorphism, and pass to the same WD class.
\end{proof}
% END GUA3

% BEGIN GUA4
\paragraph{GUA4. The required generic unitary factorization and ordering bound.}
\label{claim:GUA4}
Let $F$ be any finite extension of $\mathbb Q_p$ and let
$\rho$ be a smooth irreducible, unitarisable, generic
representation of $\operatorname{GL}_a(F)$, $a\geq1$.
There is an actual normalized-product isomorphism
\[
       \rho\simeq\mathop{\times}_{j=1}^b\delta_j\nu^{s_j},
       \qquad -\tfrac12<s_j<\tfrac12,                 \tag{GUA.7}
\]
where every $\delta_j$ is unitary square integrable and is
represented by a segment $\Delta_j$.  If two unitary
segment representations correspond to $\Delta,\Delta'$,
then
\[
            \Delta\nu^x\text{ precedes }\Delta'\nu^y
                       \quad\Longrightarrow\quad x+1\leq y.
                                                        \tag{GUA.8}
\]
\begin{proof}
BCS8 is the single checked published input for (GUA.7):
M\oe glin--Waldspurger, \emph{Le spectre r\'esiduel de
$\mathrm{GL}(n)$}, \S\S I.10--I.11, pp.~639--640.
The field here is nonarchimedean of characteristic zero,
the representation is complex, smooth, irreducible,
generic and unitarisable, and these are precisely its
hypotheses.  No rank bound or exclusion of a residue
prime occurs.  Retain its actual normalized-product
isomorphism and its strict inequalities.  The retained
square-integrable segment theorem supplies each $\Delta_j$.

For completeness, the numerical ordering bound is elementary.
Let $r,r'$ be the lengths of $\Delta,\Delta'$ and $d$ their
common cuspidal rank when the shifted segments are linked.
If $\omega$ is the central character of the first term of
$\Delta$, unitarity of its segment representation gives,
on a uniformizer,
\[
 \log|\omega(\varpi)|
          =-d\frac{r-1}{2}\log|\varpi|_F.            \tag{GUA.9}
\]
Indeed its central character is the product of the $r$
cuspidal central characters at shifts $0,\ldots,r-1$;
its norm is $1$, and $\sum_{j<r}j=r(r-1)/2$.
Precedence means that the first term of
$\Delta'\nu^y$ is the $j$th term along $\Delta\nu^x$,
with $j\geq1$, and that $r<j+r'$, because the second
segment is not contained in the first.  Applying (GUA.9)
to both segments gives
$x+j-y=(r-r')/2$.  Since $j\geq1$ and
$j+r'-r\geq1$, it follows that
$y-x=(j+j+r'-r)/2\geq1$.  This proves (GUA.8).
\end{proof}
% END GUA4

% BEGIN GUA5
\paragraph{GUA5. The local Speh parameter, with every length.}
\label{claim:GUA5}
Let $P$ be an original cuspidal representation of
$\operatorname{GL}_a(\mathbb A_E)$ over a totally complex
number field, let $\rho=P_w$ be a specified unitary local
component, and let $d\geq1$.  If the actual smooth
irreducible $\eta$ is the source's \texttt{IsSpeh}
quotient of
\[
           \rho\nu^{(d-1)/2}\times\rho\nu^{(d-3)/2}
                                  \times\cdots\times\rho\nu^{-(d-1)/2},
\]
then for every $A,h,\mathrm{rc}$ with \texttt{IsRecT},
\[
 \mathrm{rc}(\operatorname{cl}\eta)
       =\bigoplus_{k=0}^{d-1}
          \mathrm{rc}(\operatorname{cl}\rho)((d-1)/2-k).
                                                        \tag{GUA.10}
\]
\begin{proof}
BCS5 proves that the original cusp $P$ is globally generic
and that its specified local component is generic.
One can see the local implication directly: restrict a
nonzero global Whittaker functional to a finite sum of
images of the local component; one image has nonzero
restriction, furnishing the local Whittaker functional.
Apply GUA4 and write $\rho$ as the product (GUA.7).
If $\Delta_j\nu^{s_j}$ preceded
$\Delta_{j'}\nu^{s_{j'}}$, (GUA.8) would give
$s_j+1\leq s_{j'}$, impossible since both exponents
lie in $(-1/2,1/2)$.  Consequently this product is a
standard module in any ordering of these factors; its
unique quotient is $\rho$ itself.

Use induction in stages to replace the displayed Speh
induction by the induction of the factors
\[
       \delta_j\nu^{s_j+(d-1)/2-k}
                  \qquad(0\leq k<d,\ 1\leq j\leq b),
\]
ordered first by increasing $k$.  If an earlier segment
preceded a later one, GUA4 would require
\[
 s_j+(d-1)/2-k+1
                  \leq s_{j'}+(d-1)/2-k',\qquad k\leq k'.
\]
This implies $1\leq s_{j'}-s_j-(k'-k)<1$, a
contradiction.  The standard-cosocle theorem therefore
identifies every irreducible quotient of this actual
induction, including the specified $\eta$, with its
Langlands quotient.  Apply the Langlands and real-twist
laws of the given \texttt{IsRecT}.  The segment sum of
the flattened list is the sum over $k$ of the segment
sum for $\rho$, twisted by $(d-1)/2-k$.
Tensoring a segment parameter with a scalar unramified
character commutes with its special block and with finite
sums, so this is exactly (GUA.10), with its monodromy,
not just an equality of characteristic polynomials.
\end{proof}
% END GUA5

% BEGIN GUA6
\paragraph{GUA6. The total parameter identity \texttt{RecTotT}.}
\label{claim:GUA6}
Let $a_i,d_i\geq1$, $\sum_i d_i a_i=n$; let
$\rho_{i,w}$ and $\eta_{i,w}$ be the unitary cusp and
Speh local representations supplied by BCS.  Put
\[
 e_{ik}=d_i-2k-n,\qquad
 p_{w,A,\mathrm{rc}}
   =\bigoplus_i\bigoplus_{k=0}^{d_i-1}
                 \mathrm{rc}(\operatorname{cl}\rho_{i,w})(e_{ik}/2).
                                                        \tag{GUA.11}
\]
If an irreducible $\tau$ is isomorphic to the full original
normalized product $\mathop{\times}_i\eta_{i,w}$, then
for every permitted $A,h,\mathrm{rc}$,
\[
                  p_{w,A,\mathrm{rc}}
                       =\mathrm{rc}(\operatorname{algNorm}\tau).
                                                        \tag{GUA.12}
\]
\begin{proof}
GUA3 gives the parameter of the full product as the sum
of the parameters of its actual factors.  Substitute
GUA5 for each factor.  The \texttt{IsRecT} twist law
then applies to the single normalization twist
$(1-n)/2$.  Scalar twists distribute over block sums,
and
\[
        \frac{d_i-1}{2}-k+\frac{1-n}{2}
                              =\frac{d_i-2k-n}{2}.
\]
This is (GUA.11), in the same order as the original
\texttt{totPar}.  The dimension reindexing in
$\sum_i d_i a_i=n$ is a permutation of the explicitly
ordered block basis.  It cannot change the WD class.
\end{proof}
% END GUA6

\subsection{Polarized cuspidal attachment and half shifts}

% BEGIN GUA7
\paragraph{GUA7. Exactly the toral branch of \texttt{RowT}.}
\label{claim:GUA7}
Assume the original \texttt{RowT}.  Let $P$ be regular
algebraic cuspidal of rank $a\geq1$, polarized by the
original Hecke character $\chi$ of $E^+$, with
specified local components $\rho_w$.  There are $r,H$
with
\[
 (r,\,\mathrm{rc}(\operatorname{algNorm}\rho_w),\,H)
                                      \text{ satisfying \texttt{CompatG}},
\]
and
\[
   \iota H_\tau
       =\left\{\frac{a-1}{2}-z:
             z\in\operatorname{Inf}(P,\iota\tau)\right\}.
                                                        \tag{GUA.13}
\]
Here $\iota H$ means the integer multiset mapped to $\mathbb C$.
\begin{proof}
The source's \texttt{IsPolarized} first gives the sign
$\chi_v(-1)=(-1)^a$ and an isomorphism with the conjugate
inverse-transpose twist.  The retained
Gelfand--Kazhdan duality identifies inverse transpose
with the smooth contragredient.  Pair the resulting
contragredient with $P$ by evaluation.  This is
nondegenerate in each variable; finite adelic operators
act contragrediently, compact operators act by the
corresponding inverse, and differentiating gives the
sum-zero identity for Lie operators.  Thus it is
the literal \texttt{PolarizedPair} required by
\texttt{RowT}, including its sign.

Apply \texttt{RowT}, which is \texttt{Row0b(IsRecT)},
and retain its actual $r$, $H$ and all fields of
\texttt{PropsB}.  Continuity, semisimplicity and
(GUA.13) are three of those fields.  For $w\nmid\ell$
take the actual monodromy parameter $\operatorname{wdOf}(r)$.
The \texttt{away} field gives \texttt{AwayB} against
the parameter of every permitted $\mathrm{rc}$.  Its
\texttt{tor} clause says these are equal when
\texttt{Tor} and $N=0$ hold, proving (GUA.3).

For every $O,B_w,\texttt{Induces}$ use the \texttt{above}
field with the original tied period model
$\texttt{Loc.PP}(\ell)$.  Its \texttt{nice} alternative
contains \texttt{Tor} as a separate disjunct, followed
by $N=0$.  That branch gives de Rhamness, potential
semistability, every prescribed filtration-jump statement,
and \texttt{AwayB} for the actual $\operatorname{wdAt}$.
Applying the latter's toral clause proves (GUA.4).
If the parameters are good for all correspondences,
the retained \texttt{RecTExists} supplies one at each
$A,h$ to apply this branch; its jump conclusion does
not depend on which one was chosen.  This proves
\texttt{HasHTAt}.  No use was made of the other
\texttt{nice} branch, which asks for $\mathbb Q_p$
and a rank-dependent lower bound on $p$.  In particular
this proof imposes no such bound and no exclusion of $\ell=3$.
\end{proof}
% END GUA7

% BEGIN GUA8
\paragraph{GUA8. The half-shift construction \texttt{HalfG}.}
\label{claim:GUA8}
Let $P$ be a conjugate-self-dual cusp of rank $a\geq1$,
with actual local components $\rho_w$ and infinity
multisets $m_\phi$ satisfying $P.\texttt{HasInfChar}(\phi,m_\phi)$.
Suppose $e\in\mathbb Z$, every $m_\phi$ has distinct
entries, and $z+e/2\in\mathbb Z$ for every entry.
There are $r,H$ satisfying \texttt{CompatG} with
\[
       p_{w,A,\mathrm{rc}}
                       =\mathrm{rc}(\operatorname{cl}\rho_w)(e/2),
       \qquad
       \iota H_\tau=\{-z-e/2:z\in m_{\iota\tau}\}.
                                                        \tag{GUA.14}
\]
\begin{proof}
Set $s=(e+a-1)/2$ and form the \emph{actual}
automorphic twist $P_s=P\otimes|\det|_{\mathbb A_E}^{s}$.
Multiplication of functions by the determinant character
is an invertible linear map onto this cusp subspace;
constant terms remain zero, and right translation is
twisted by $|\det|^s$.  Its restriction at $w$ therefore
carries the given local intertwiner for $\rho_w$ to
an intertwiner for $\rho_w\nu^s$.  At infinity the
infinitesimal character is $m_\phi+s$.  These entries
are distinct and lie in $\mathbb Z+(a-1)/2$, so this
is regular algebraic, with the source normalization
also in rank $1$.

Write $\xi=|\cdot|_{\mathbb A_E}^{s}$ and let
$\operatorname{res}_{E^+}\xi$ be its restriction to
the ideles of $E^+$.  Its norm pullback is
$\xi\xi^c=\xi^2$.  Let $\delta_{E/E^+}$ be the
quadratic Hecke character.  It is trivial on norms
and has value $-1$ at $-1$ at every real place.
Consequently
\[
        \chi=(\operatorname{res}_{E^+}\xi)
                                  \delta_{E/E^+}^{a}
\]
has $\chi_v(-1)=(-1)^a$ and norm pullback $\xi^2$.
Transport the given conjugate-self-duality intertwiner
of $P$ across multiplication by $\xi$.  The identity
$\xi(\xi\xi^c)^{-1}=(\xi^c)^{-1}$ verifies, for
each finite, compact and Lie operator, precisely the
conjugate-dual twist identity for $P_s$ and $\chi$.
Thus $P_s$ has the polarization used in GUA7.

Apply GUA7.  For each individual permitted $\mathrm{rc}$,
its real-twist law gives
\[
 \mathrm{rc}(\operatorname{algNorm}(\rho_w\nu^s))
  =\mathrm{rc}(\operatorname{cl}\rho_w)
                       \bigl(s+(1-a)/2\bigr)
  =\mathrm{rc}(\operatorname{cl}\rho_w)(e/2).
\]
Substituting this equality preserves every compatibility
and conditional-goodness quantifier.  Formula (GUA.13)
becomes $(a-1)/2-(z+s)=-z-e/2$, proving (GUA.14).
The representation was produced by the original
attachment theorem applied to $P_s$.  This construction
does not require a square root of the cyclotomic
character as a Galois character.
\end{proof}
% END GUA8

% BEGIN GUA9
\paragraph{GUA9. Infinity weights and the exact Speh strings.}
\label{claim:GUA9}
At every embedding $\phi:E\to\mathbb C$ there is a
dominant integral weight $\lambda$ with
$\texttt{HasWeightAt}(\sigma,\phi,\lambda)$.  Suppose
the BCS infinity identity is
\[
 \coprod_i\coprod_{k=0}^{d_i-1}
       \{z+(d_i-1)/2-k:z\in m_{i,\phi}\}
     =\operatorname{infOfWeight}(\lambda)
     =\{\lambda_j+(n-1)/2-j:0\leq j<n\}.           \tag{GUA.15}
\]
Then each $m_{i,\phi}$ has distinct entries, and for
each $k<d_i$ all $z+e_{ik}/2$ are integers.  Moreover
\[
 \coprod_{i,k}\{-z-e_{ik}/2:z\in m_{i,\phi}\}
   =\{(n-1)/2-z:z\in\operatorname{infOfWeight}(\lambda)\}.
                                                        \tag{GUA.16}
\]
\begin{proof}
We first justify the existence of a weight on the actual
automorphic function space.  The infinity orbit of a
nonzero function spans a nonzero finite-dimensional space
by \texttt{IsAutRepU}.  Evaluation at group elements
separates this space, so finitely many evaluations form
a basis of its dual.  Continuity of the functions and
of right translation therefore proves continuity of
its matrix representation on the compact factor
$U(n)$ at $\phi$.  Averaging an inner product splits
off a finite-dimensional irreducible continuous module $S$.

Here the compact highest-weight classification can be
obtained from the already proved HP4 without introducing
a new automorphic theorem.  The restrictions of regular
functions on $\operatorname{GL}_n(\mathbb C)$ form a
self-conjugate algebra separating points of $U(n)$:
on $U(n)$, $\overline{g}_{ij}=(g^{-1})_{ji}$.
Stone--Weierstrass makes them dense in the continuous
functions.  A nonzero matrix coefficient of $S$ thus
has nonzero Haar pairing with some such regular function.
The finite-dimensional right-translation space generated
by that function is a rational $\operatorname{GL}_n$
module.  The Haar pairing, or its matrix-coefficient
projector, gives a nonzero $U(n)$-map from this space
to $S$; invariant orthogonal complements split it.
The group $U(n)$ is Zariski dense in
$\operatorname{GL}_n(\mathbb C)$: the Cayley map from
skew-Hermitian matrices has real coordinates whose
complexification is the full matrix space, and a complex
polynomial vanishing on all those real coordinates
vanishes identically.  Hence the resulting irreducible
$U(n)$ summand is an irreducible rational
$\operatorname{GL}_n$ submodule.  HP4 identifies it
with the literal polynomial module
$\texttt{sqRep}(\lambda)$ for a dominant integral
$\lambda$.  Its nonzero intertwiner into the original
infinity orbit is exactly the definition of
\texttt{HasWeightAt}.  This argument also preserves
the chosen embedding, including its conjugate orientation.

The right side of (GUA.15) has distinct entries:
if $j<j'$, dominance gives
$\lambda_j-j>\lambda_{j'}-j'$.  A submultiset of a
multiset with distinct entries has distinct entries,
and translation is injective; since $d_i\geq1$ this
proves the first assertion.  If an entry on the left
equals the $j$th entry on the right, subtract
$(n-1)/2$ to obtain
\[
          z+(d_i-2k-n)/2=\lambda_j-j\in\mathbb Z.
\]
Apply $u\mapsto(n-1)/2-u$ to both sides of
(GUA.15) to obtain (GUA.16), with multiplicities.
The argument applies to every weight witness in
(GUA.15), not only the one used to prove existence.
\end{proof}
% END GUA9

\subsection{One representation and its actual period modules}

% BEGIN GUA10
\paragraph{GUA10. Finite direct sums preserve the full compatibility record.}
\label{claim:GUA10}
Given finitely many triples $(r_j,p_j,H_j)$ satisfying
\texttt{CompatG}, their actual block sum satisfies
\texttt{CompatG} with parameter $\bigoplus_jp_j$
and jump multiset $\coprod_jH_j$.  If there is at least
one block and all block ranks are positive, irreducibility
of the sum forces exactly one block, which is irreducible.
This proves \texttt{SumHTAt}, \texttt{SumTwoG} and
the full dimension-reindexed \texttt{SumFamG}.
\begin{proof}
Use the ordered direct sum of the given coordinate spaces;
on it define $R(g)=\operatorname{diag}(r_j(g))$.
Entries and inverse entries are continuous functions of
those of the constituents.  It is semisimple because a
finite direct sum of semisimple modules is semisimple.
For the remaining assertions, the coordinate inclusions
$u_j$ and projections $v_j$ satisfy
\[
 v_j u_k=\delta_{jk},\qquad \sum_j u_jv_j=1,
\]
and intertwine the actual Galois actions.

Here are the needed period compatibilities explicitly.
Choose a common finite coefficient extension for the
finitely many models in any period assertion.  It is
obtained by adjoining their coefficient fields; extension
of a finite coefficient field commutes with fixed spaces,
since a basis of the larger field expands the fixed
equations coefficient by coefficient and the coefficient
Galois action is trivial.  Tensor the displayed maps with
the identity of the specified period ring.  They commute
with the \emph{diagonal} Galois action: both the
representation coordinate and the period coordinate
transform on each side.  Thus the same two identities
hold on invariants, and give
\[
 D_{\rm dR}(R)=\bigoplus_jD_{\rm dR}(r_j),\qquad
 D_{\rm pst}(R)=\bigoplus_jD_{\rm pst}(r_j).        \tag{GUA.17}
\]
For $D_{\rm pst}$ this is first true at every common
finite extension of the local field.  In the union over
such extensions, finitely many components of an element
are all defined at one common extension, so (GUA.17)
continues to hold.  The same proof applies to the
maximal-unramified scalar version in \texttt{IsPst}.
It works in the source's actual tied period model,
using its own rings and actions.

The maps also commute with $\varphi$ and $N$, which
act on the period coordinate, and preserve each
filtration level.  The de Rham and potentially
semistable dimension equalities therefore add.
Every graded filtration dimension adds at each
specified coefficient embedding, giving multiset union
of the jumps for \emph{all} embeddings in
\texttt{HasHTAt}.  The WD construction from
$D_{\rm pst}$, including Frobenius semisimplification,
is consequently the block sum of the constituent
constructions.  Semisimplification commutes with a
finite block sum, since the Jordan decomposition of a
block-diagonal matrix is block diagonal.

Away from $\ell$, choose a common finite coefficient
field (ADK1) and a common open inertia subgroup for
the finitely many monodromy descriptions.  On it the
representations are $\exp(t_\ell(g)N_j)$.
Their block sum is $\exp(t_\ell(g)\bigoplus_jN_j)$;
removing this exponential and choosing the same
Frobenius gives the sum of their Weil representations.
This proves the asserted sum identity for the actual
\texttt{IsWDOf} parameters.

To check the conditional statements, GUA1 says goodness
of the total parameter implies goodness of every
summand for that same $A,h,\mathrm{rc}$.  The preceding
period and WD identities therefore prove (GUA.3)--(GUA.4).
If the total is good for every such triple, each constituent
satisfies the premise of its own \texttt{HasHTAt};
(GUA.17) proves the total jump assertion.  This keeps
both universal quantifiers in their original order.

Finally, if two positive blocks occur, the first block's
coordinate space is a nonzero proper stable subspace.
If exactly one block occurs but is reducible, its
nonzero proper stable subspace is also one for the total.
Reindexing the block coordinates by the specified
equality of their rank sum with $n$ is a fixed linear
permutation and transports all these statements.
\end{proof}
% END GUA10

% BEGIN GUA11
\paragraph{GUA11. Construction of the one global $R$ and $H$.}
\label{claim:GUA11}
Choose the BCS witnesses for the original $\sigma$:
finitely many unitary conjugate-self-dual cusps $P_i$
of ranks $a_i\geq1$, lengths $d_i\geq1$, actual
local cusp and Speh components $\rho_{i,w},\eta_{i,w}$,
and actual infinity multisets $m_{i,\phi}$, with
$\sum_i a_i d_i=n$.  There are a continuous semisimple
representation
\[
           R:G_E\longrightarrow\operatorname{GL}_n(
                                  \overline{\mathbb Q}_\ell)
\]
and integer multisets $H_\tau$ satisfying
\texttt{CompatG} for (GUA.11).  Every infinity-weight
witness $\lambda$ at $\iota\tau$ satisfies
\[
 \iota H_\tau
       =\{(n-1)/2-z:z\in\operatorname{infOfWeight}(\lambda)\}.
                                                        \tag{GUA.18}
\]
If $R$ is irreducible, then there is just one $i$ and
its $d_i$ is $1$.
\begin{proof}
The number of cusps is positive since otherwise their
positive rank sum would be $0=n$.  For each pair $(i,k)$,
$0\leq k<d_i$, GUA9 proves that $m_{i,\phi}$ and
$e=e_{ik}$ satisfy GUA8 at every $\phi$; existence
of an infinity weight was proved for each embedding
before making this inference.  Apply GUA8 to this
particular cusp, its entire family of specified local
components, and this one integer $e_{ik}$.  It gives
an actual representation $r_{ik}$ of rank $a_i$ and
multisets $H_{ik,\tau}$, simultaneously at all places
and all correspondences.

Take the ordered block sum of these representations,
first over $k$, then over $i$, and transport the basis
by $\sum_i d_i a_i=n$.  Define
$H_\tau=\coprod_{i,k}H_{ik,\tau}$.  GUA10 proves
the entire compatibility record for this one $R$.
The two successive sums give exactly the two lists in
the source's \texttt{totPar}, with no change of map
between different local conditions.  Substitution of
(GUA.14) and then (GUA.16) proves (GUA.18).
All blocks have positive rank.  Hence GUA10 first
forces one outer block, and then one inner block,
under irreducibility.  These are precisely $c=1$
and $d_0=1$; they were not premises of the construction.
\end{proof}
% END GUA11

\subsection{All unramified places and every split local component}

% BEGIN GUA12
\paragraph{GUA12. Spherical normalization for every \texttt{IsRecT}.}
\label{claim:GUA12}
If a smooth irreducible $\rho$ of $\operatorname{GL}_a(E_w)$
has normalized Satake multiset $S=\{\alpha_1,\ldots,\alpha_a\}$,
then for every permitted $A,h,\mathrm{rc}$ the class
$\mathrm{rc}(\operatorname{cl}\rho)$ is unramified,
has $N=0$, and has arithmetic Frobenius eigenvalues
$\alpha_1^{-1},\ldots,\alpha_a^{-1}$.
\begin{proof}
We spell out the needed spherical specialization of the
retained local classification, rather than assuming that
a quotient of a spherical induction is automatically
spherical.  BCS14 proves exactness of maximal-compact
fixed vectors and their tensor formula for induction.
BCS15 proves that the spherical Hecke algebra is generated
by $T_1,\ldots,T_a,T_a^{-1}$, computes its normalized
characters, and proves that an irreducible spherical
fixed space is one dimensional.

We also need uniqueness for a given Hecke character,
which has the following algebraic proof.  Let
$X=\mathrm{c\!\operatorname{-}Ind}_K^G\mathbf1$ and let
$\mathcal H=\operatorname{End}_G(X)$, with the right
module convention on $X$.  For a Hecke character $\lambda$,
$M=X\otimes_{\mathcal H}\mathbb C_\lambda$ is generated
by its single $K$-fixed line.  Averaging by the idempotent
$e_K$ shows that $M^K=\mathbb C_\lambda$.  Every proper
submodule has zero $K$-fixed space, since otherwise it
contains this generating line.  The sum of all proper
submodules still has zero fixed space (averaging a
finite sum is the sum of the averages), hence is proper.
Its quotient is the unique irreducible spherical quotient.
Every irreducible spherical representation with character
$\lambda$ is a quotient of $M$ by the compact-induction
adjunction.  Consequently two such representations are
isomorphic.  The unramified principal series of the
characters $\xi_j(\varpi)=\alpha_j$ has this character
by BCS15; its fixed-generated submodule has this unique
quotient.  Thus $\rho$ is a subquotient of this actual
principal series.

Choose the Langlands multisegment of $\rho$ by the
retained exhaustion theorem.  Its standard-module
surjection and BCS14 force each segment representation
to be spherical.  A spherical segment representation is
a quotient of the induction of its cuspidal terms, so
BCS14 makes each such term spherical.  A spherical
supercuspidal is an unramified character: the preceding
principal-series realization gives its support as a
list of unramified rank-one characters, whereas its
supercuspidal support is the singleton consisting of
itself.  The retained support uniqueness theorem forces
rank one and that character.

It remains to rule out a spherical segment of length
$t>1$.  Write its first term as the unramified character
$\mu$, so its induction has characters
$\mu,\mu\nu,\ldots,\mu\nu^{t-1}$.  The determinant
character $\eta\circ\det$, where
$\eta=\mu\nu^{(t-1)/2}$, embeds in that induction:
the function $g\mapsto\eta(\det g)$ satisfies its
normalized covariance exactly.  Both this character
and the purported spherical segment have the same
spherical Hecke character, by BCS14--BCS15.  The
uniqueness just proved identifies them.  But the
Langlands multisegment of the determinant character
consists of the $t$ singleton characters in descending
order.  To verify this last fact, embed the inverse
determinant character in its ascending induction by
the same formula, take smooth contragredients, and use
the retained normalized-induction duality; this yields
a surjection of the descending singleton induction
onto $\eta\circ\det$.  That list has no earlier
segment preceding a later one, so the standard-cosocle
theorem gives its Langlands multisegment.  Uniqueness
compares its $t$ segments with the single length-$t$
segment, forcing $t=1$.

Thus all Langlands segments of $\rho$ are unramified
singletons.  BCS15 on their standard induction and
BCS14 on its quotient identify their uniformizer values
with $S$.  This establishes the required singleton
Langlands description, including its ordering and
normalization, for the given $\rho$.

The Langlands and rank-one laws of the given
\texttt{IsRecT} therefore give
\[
          \mathrm{rc}(\operatorname{cl}\rho)
                     =\bigoplus_{j=1}^a(\xi_j\circ a_{\mathrm{rc}},0).
\]
The degree convention in $\texttt{IsLocalArtin}$
gives valuation $-1$ on arithmetic Frobenius, so the
unramified characters have values $\alpha_j^{-1}$
there.  Their inertia action is $1$.  This proves
the statement at the level of the actual WD class.
For rank zero the sum is empty.  In particular the
calculation is independent of any choices in
$\mathrm{rc}$, even though its full predicate is
stronger than the usual unramified correspondence.
\end{proof}
% END GUA12

% BEGIN GUA13
\paragraph{GUA13. Local Satake data on the original cusp subspace.}
\label{claim:GUA13}
For each specified cusp component $\rho_{i,w}$, its
local Satake multiset is also the multiset in the literal
$P_i.\texttt{HasSat}(w,-)$, and conversely.  Such
multisets exist for all $i$ at a cofinite set of places.
They also exist for every $w$ satisfying the original
$\texttt{IsUnramU}(\sigma,w)$.
\begin{proof}
The restriction of an irreducible cusp representation to
one local factor is isotypic of any specified local
component.  Indeed the complementary adelic factors
and the infinity operators commute with this factor.
Translates of one nonzero local intertwiner are again
images of that same local representation, and their
sum is a nonzero subspace stable under every global
operator.  Irreducibility makes their sum the whole
cusp space.  Equivalently this is the elementary
isotypy argument of ITA10.

Let $K=\operatorname{GL}_{a_i}(\mathcal O_{E_w})$.
A local $K$-fixed eigenvector embeds into the global
cusp space through a nonzero intertwiner, which is
injective by local irreducibility.  Conversely, write
a global $K$-fixed vector as a finite sum of images
of local vectors, and average those vectors over $K$.
The average is a finite sum on each smooth vector;
the sum of the images is still the original nonzero
vector, so at least one local average is nonzero.
The same argument shows that every global fixed
vector has the local Hecke eigencharacter.  Uniqueness
of the spherical Satake multiset gives the asserted
equivalence, including the source requirement on all
nonzero fixed vectors.

Choose a nonzero smooth finite-adelic vector in each
of the finitely many $P_i$.  Its open stabilizer contains
the standard maximal compact at all but finitely many
places, by the restricted-product topology.  The
previous paragraph makes all the $\rho_{i,w}$
spherical there; the spherical Hecke classification
then supplies their Satake multisets.  For every
$\texttt{IsUnramU}(\sigma,w)$, BCS19 instead supplies spherical $\rho_{i,w}$
at that very place, so the same argument gives the
last assertion without discarding an extra finite set.
\end{proof}
% END GUA13

% BEGIN GUA14
\paragraph{GUA14. \texttt{FrobTotT} and the complete unramified assertion.}
\label{claim:GUA14}
Let $w\nmid\ell$ and suppose the specified
$\rho_{i,w}$ have Satake multisets $S_{i,w}$.
Define, with multiplicities,
\[
 B_w=\coprod_i\coprod_{k=0}^{d_i-1}
       \{\alpha q_w^{k-(d_i-1)/2}:\alpha\in S_{i,w}\}.
                                                        \tag{GUA.19}
\]
For any continuous $R$ satisfying
$\texttt{ParAtT}(\iota,R,w,p_w)$ with (GUA.11),
all inertia groups over $w$ act trivially, and for
every valuation $A$ over $w$ and every arithmetic
Frobenius $g$ at $A$,
\[
 \iota\det(X-R(g^{-1}))
   =\prod_{\beta\in B_w}(X-q_w^{(n-1)/2}\beta).
                                                        \tag{GUA.20}
\]
This is the entire \texttt{FrobEig} assertion, hence
\texttt{FrobTotT}.  The $R$ of GUA11 is consequently
unramified at \emph{every} $w\nmid\ell$ with
$\texttt{IsUnramU}(\sigma,w)$.
\begin{proof}
Fix any $A,h$ and use the retained \texttt{RecTExists}
to choose one permissible $\mathrm{rc}$ there.  GUA12
and (GUA.11) make $p_{w,A,\mathrm{rc}}$ unramified,
with arithmetic eigenvalues
$\alpha^{-1}q_w^{e_{ik}/2}$ and zero monodromy.
It is good by GUA1.  Thus the conditional equality
in \texttt{ParAtT} identifies it with the actual
Weil--Deligne parameter of $R$ at $A$.

This equality implies trivial \emph{original} inertia
action.  In the literal \texttt{IsWDOf} witness, the
zero monodromy operator makes $R$ trivial on an open
subgroup of inertia.  Compactness of inertia then
makes its image finite, so each inertia matrix is
semisimple in characteristic zero.  The original
\texttt{IsFrobSS} witness commutes with that matrix
and differs from it by a nilpotent matrix; since
both commuting matrices are semisimple, simultaneous
diagonalization makes their nilpotent difference zero.
Thus the original inertia matrix equals the inertia
matrix of the unramified Weil--Deligne parameter,
which is $1$.  For general Weil elements the same
commuting nilpotent difference preserves the
characteristic polynomial.
Inverting the arithmetic eigenvalues gives
\[
 \alpha q_w^{-e_{ik}/2}
    =q_w^{(n-1)/2}\alpha q_w^{k-(d_i-1)/2},
\]
which proves (GUA.20).  Trivial inertia makes it true
for every arithmetic Frobenius lift, and $A,h$ were
arbitrary.  The proof uses only the hypotheses of
the literal \texttt{FrobTotT}, so it is not confined
to the representation constructed in GUA11.
For that representation use GUA13 at every
$\texttt{IsUnramU}$ place.  No cofinite exception
is introduced in this last conclusion.
\end{proof}
% END GUA14

% BEGIN GUA15
\paragraph{GUA15. Every split local component and its conditional jumps.}
\label{claim:GUA15}
For every split $w$ and every original local component
$\tau$ of $\sigma$ along the specified map
$\texttt{cLoc}\circ\texttt{locU}$, the representation
of GUA11 satisfies
\[
       \texttt{LocGlobT}(\iota,R,w,\tau),\qquad
       \texttt{PlainTor}(\operatorname{cl}\tau)
                        \Longrightarrow\texttt{HasHTAt}(R,H,w).
                                                        \tag{GUA.21}
\]
\begin{proof}
At a split place the local unitary group is the
specified copy of $\operatorname{GL}_n(E_w)$; in
the paired coordinate the unitary equation determines
the other matrix as the conjugate inverse transpose.
The adelic group is a product of this factor and
its centralizer.  If an irreducible global module
contains one local irreducible image $V$, its global
translates are again images of $V$: remove the local
factor of the translating element, and the remaining
factor commutes with the local group.  The images
span the whole irreducible module.  ITA10, or a
coordinate projection from a finite sum of these
simple images, now shows that any second nonzero
local irreducible image is isomorphic to $V$.
This proves local-component uniqueness on the actual
function subspace, with no tensor-product realization
of that subspace assumed.

BCS supplies the full local normalized product
$\mathop{\times}_i\eta_{i,w}$ as a local component
of $\sigma$.  It is irreducible by the retained
unitary-product theorem (each factor is unitary),
or directly by the irreducibility field of its
$\texttt{IsCompAlong}$ witness.  The uniqueness just
proved identifies the given $\tau$ with that exact
product.  GUA6 identifies (GUA.11) with
$\mathrm{rc}(\operatorname{algNorm}\tau)$ for
every $A,h,\mathrm{rc}$, so GUA11's full
\texttt{ParAtT} is exactly \texttt{LocGlobT}.
If $\tau$ is \texttt{PlainTor}, GUA2 makes this
parameter good for every one of those triples.
The \texttt{ht} field of GUA11's \texttt{CompatG}
therefore gives \texttt{HasHTAt}, with all its
original coefficient-ring, bridge and embedding quantifiers.
\end{proof}
% END GUA15

\subsection{Polarization and the irreducible case}

% BEGIN GUA16
\paragraph{GUA16. The required pairing for every complex conjugation.}
\label{claim:GUA16}
For every $c\in G_{E^+}$ satisfying the source's
\texttt{IsConjEl}, the $R$ of GUA11 admits an
invertible matrix $P$ with
\[
       R(s)^t P R(csc^{-1})=\epsilon(s)^{1-n}P
                                      \qquad(s\in G_E). 
                                                        \tag{GUA.22}
\]
Equivalently this is the source's \texttt{IsPairing}
for every pair $s,s'$ whose restrictions satisfy
$s'=csc^{-1}$ in $G_{E^+}$.
\begin{proof}
At a cofinite set of places, all the cusp Satake
multisets in GUA13 exist at both $w$ and $cw$, and
both places are away from $\ell$.  Complex
conjugation preserves their residue norms.  The
conjugate-self-duality intertwiner of $P_i$ takes
its local spherical Hecke character at $cw$ to the
inverse Satake character at $w$.  To check the
normalization, inverse transpose sends the diagonal
torus character to its inverse; the half-modulus
normalization of the Satake transform is invariant
under reversing the roots.  Thus
\[
                 S_{i,cw}=\{\alpha^{-1}:\alpha\in S_{i,w}\}.
\]
Reverse $k$ to $d_i-1-k$ in (GUA.19).  Its exponent
changes sign, so $B_{cw}=B_w^{-1}$, including every
Speh length.  Let $C_w=q_w^{(n-1)/2}$.
By GUA14, $R^c$ has geometric Frobenius eigenvalues
$C_wB_{cw}$ at $w$.  On the same lift
$R^\vee\otimes\epsilon^{1-n}$ has eigenvalues
\[
         q_w^{n-1}(C_wB_w)^{-1}
                  =C_wB_{cw}.                         \tag{GUA.23}
\]
Here $\epsilon$ on geometric Frobenius is $q_w^{-1}$,
so the exponent $1-n$ in (GUA.23) is forced.

Both representations are continuous and semisimple.
Conjugating the group preserves these properties;
duality preserves semisimplicity by annihilators
of submodules, and a character twist preserves
the stable-subspace lattice.  Chebotarev makes
the union of the Frobenius conjugacy classes
outside the finite exceptional set dense in $G_E$.
Their characteristic polynomials agree by (GUA.23).
Each coefficient is a continuous function into
the Hausdorff field $\overline{\mathbb Q}_\ell$,
so equality holds on all of $G_E$.  The
Brauer--Nesbitt argument then gives an isomorphism
of the two semisimple representations: in the
finite-dimensional algebra spanned by their
combined image, simple-module characters separate
the semisimple factors, and equality of traces
gives their multiplicities.  This needs no purity
statement or bound on Satake absolute values.

Choose the matrix $P$ of an intertwiner
$R^c\longrightarrow R^\vee\epsilon^{1-n}$.
Its identity is
$P R(csc^{-1})=\epsilon(s)^{1-n}R(s)^{-t}P$;
multiplication by $R(s)^t$ gives (GUA.22).
Restriction of automorphisms is injective here,
so the formulation with arbitrary $s'$ is the
same identity.  The construction applies to each
specified complex conjugation $c$.  No symmetric
or alternating sign for $P$ is required by this endpoint.
\end{proof}
% END GUA16

% BEGIN GUA17
\paragraph{GUA17. Irreducibility produces the actual cuspidal representation.}
\label{claim:GUA17}
If $R$ is irreducible, there is an original cusp
$\Pi$ of $\operatorname{GL}_n(\mathbb A_E)$ that is
regular algebraic and satisfies the source's
\texttt{DualIsConj} and \texttt{AttC}.  It has
all the infinity weights of $\sigma$, has a
Satake multiset at every
$\texttt{IsUnramU}(\sigma,w)$, and every one of
its local components at a split $w$ is a component
of $\sigma$ along the original local map.
\begin{proof}
GUA11 gives $c=1$, $d_0=1$ and consequently $a_0=n$.
Take the actual BCS cusp $\Pi=P_0$; no automorphy
of a new representation is inferred.  BCS's
infinity equality now reads
$m_{0,\phi}=\operatorname{infOfWeight}(\lambda)$
for every $\texttt{HasWeightAt}$ witness.  Its
\texttt{HasInfChar} field gives the desired exact
infinity character.  At each embedding, GUA9
supplies such a weight; its entries are distinct
and lie in $\mathbb Z+(n-1)/2$.  This proves
regular algebraicity, also for $n=1$.

The given conjugate-self-duality of $P_0$, combined
with the retained Gelfand--Kazhdan duality and
evaluation pairing as in GUA7, gives
\texttt{DualIsConj}, including the finite, compact
and Lie identities in its definition.  At a
cofinite set of places GUA13 supplies its global
Satake multiset $S_{0,w}$.  For length one (GUA.19)
is exactly $S_{0,w}$, and GUA14 supplies
$\texttt{FrobEig}(\iota,R,w,
\operatorname{algSat}(n,q_w,S_{0,w}))$.
Together with semisimplicity this is precisely
\texttt{AttC}, not just equality of unspecified
Euler factors.  The every-unramified-place
assertion follows directly from the last clause
of GUA13, without imposing $w\nmid\ell$.

For length one the defining Speh induction is
$\rho_{0,w}$ itself, so its irreducible quotient
$\eta_{0,w}$ is isomorphic to $\rho_{0,w}$.
The one-factor normalized product is the same
representation.  The BCS split component map
therefore makes $\rho_{0,w}$ a component of
$\sigma$.  Any other local component of $P_0$
is isomorphic to it: apply the isotypy argument
of GUA13 to the cusp space.  Transporting the
actual nonzero local intertwiner along this
isomorphism proves the requested component
assertion for \emph{every} local component of $\Pi$.
\end{proof}
% END GUA17

% BEGIN GUA18
\paragraph{GUA18. The complete \texttt{GalUnitaryT} and \texttt{GuWrap} endpoint.}
\label{claim:GUA18}
Under the original hypotheses stated at the beginning,
BCS and the retained original \texttt{RowT} supply
one pair $(R,H)$ with every field of the literal
\texttt{GalUnitaryT}.  With the original
$\texttt{HypothesesH}'$ literature record, this is
the original \texttt{GuWrap.galUnitaryT\_of\_printsH}.
\begin{proof}
Choose the actual witnesses of BCS18--BCS19 and construct
$R,H$ as in GUA11.  The following table identifies
each conjunct, including its quantifiers; each
proof concerns this same $R,H$.
\begin{center}
\begin{tabular}{p{.69\linewidth}p{.19\linewidth}}
\hline
Original conclusion & Proof \\
\hline
Continuity and semisimplicity on $G_E$ & GUA11 \\
Every $w\nmid\ell$ with $\texttt{IsUnramU}$,
every valuation above it, and every inertia element:
$R(s)=1$ & GUA13--14 \\
Every $\tau$ and every infinity-weight witness:
the exact multiset in (GUA.18) & GUA9, GUA11 \\
Every split place and every specified local component:
\texttt{LocGlobT}, with every $A,h,\mathrm{rc},O,B_w$
and \texttt{Induces}; its \texttt{PlainTor}
implication to \texttt{HasHTAt} & GUA15 \\
Every original complex conjugation:
\texttt{IsPairing} with $\epsilon^{1-n}$ & GUA16 \\
If $R$ is irreducible: an actual regular algebraic
conjugate-self-dual cusp, \texttt{AttC}, all infinity
weights, all $\texttt{IsUnramU}$ Satake places and
every split local component & GUA17 \\
\hline
\end{tabular}
\end{center}
The hypotheses of BCS include exactly finite-place
unramifiedness of $E/E^+$ and the original
\texttt{IsAutRepU}; its discrete base-change input
is the retained $\texttt{litN.aut.bc}$.  The
polarized attachment is the original
$\texttt{litN.aut.gal}$, whose source wrapper
$\texttt{rowOf}'$ is the identity on \texttt{RowT}.
GUA7 used its unrestricted toral branch.  All
local arguments used only the retained classical
foundations and the checked BCS local unitary
input, and kept the full \texttt{IsRecT}
predicate.  Hence substitution of these original
inputs gives the claimed \texttt{GuWrap} theorem.
No step removes its foundational hypotheses, changes
$\ell\ne2$ to a larger-prime condition, or imposes
$d_i=1$ without the stated irreducibility premise.
\end{proof}
% END GUA18

\subsection{Scope and source boundary}

The new proofs establish the local parameter, normalization,
period-sum, block-construction, pairing and cuspidal-consumer
arguments.  The substantial inherited attachment theorem is
the full \texttt{RowT}/\texttt{PropsB} record; it has not
been replaced by a theorem-name-only assertion about Galois
representations.  Likewise BCS retains the original
\texttt{BaseChangeIsobaricN} foundation.  The existence of
full \texttt{IsRecT} correspondences and the actual tied
period model remain explicit local foundations.  The
source-pinned metadata records their precise statements
separately from the numbered-claim dependency graph.
This chapter does not assert the independent
\texttt{endShapeT} strict-norm-ratio endpoint or purity;
those conclusions have different hypotheses and consumers.
